\documentclass[12pt]{article}
\usepackage[utf8]{inputenc}
\usepackage[T1]{fontenc}
\usepackage[margin=1in]{geometry}
\usepackage{setspace}
\usepackage{xcolor}
\usepackage[colorlinks=true, urlcolor=blue!60!black, linkcolor=blue!60!black]{hyperref}
\setlength{\parindent}{0em}
\setlength{\parskip}{0.6em}
\renewcommand{\baselinestretch}{1.0}
\begin{document}
\begin{center}
{\LARGE \textit{NAACP v. Alabama ex rel. Patterson}, 357 U.S. 449 (1958)\par}
\vspace{0.5em}
{\large Case Brief --- Two-Page Edition\par}
\end{center}

\section{Memory Jogger}
Compelling the NAACP to disclose its rank-and-file members, exposing them to reprisal without a compelling, closely related state interest, violates the freedom of association protected by the Fourteenth Amendment (\href{https://supreme.justia.com/cases/federal/us/357/449/}{Justia}).

\section{Detailed Case Facts}
Alabama sued in 1956 to oust the NAACP for failing to register as a foreign corporation, citing its civil rights work. The NAACP offered to qualify and produced its charter, officers, records, and membership totals, but refused to reveal members' names.

\section{Procedural History}
Held in contempt and fined \$100,000, the NAACP was barred from a merits hearing; the Alabama Supreme Court refused review on procedural grounds.

\section{Judicial Votes}
Reversed and remanded, 9--0. Justice Harlan wrote for a unanimous Court.

\section{Holding}
Alabama's novel procedural bar was not an adequate state ground, and the NAACP could assert its members' rights. Freedom of association is part of Fourteenth Amendment ``liberty,'' and compelled disclosure of members was unjustified.

\section{Analysis of Opinions}
\textbf{Harlan:} Members could not sue in their own names without losing anonymity. State action curtailing association faces ``the closest scrutiny''; past disclosure brought ``economic reprisal, loss of employment, threat of physical coercion.'' Private reprisals counted because state compulsion triggered them. The names had no ``substantial bearing'' on whether the NAACP did business in Alabama, which it admitted. \textit{Bryant v. Zimmerman} (Klan rosters) was distinguishable.

\section{Examples: Future Applications}
\textbf{Same-Side:} A requirement that teachers list every association fails (\textit{Shelton v. Tucker}). \textbf{Opposite-Side:} Public reporting of federal campaign contributions survives (\textit{Buckley}). \textbf{Fence-Sitter:} A confidential donor-filing law aimed at foreign influence pits \textit{Bonta} against \textit{Buckley}.

\section{Critique}
Harlan derived an unenumerated right of association later invoked in \textit{Griswold}. Critics like Justice Sotomayor warn that extending it to routine donor disclosure shields ``dark money.''

\section{Key Quotations}
\begin{quote}
``It is beyond debate that freedom to engage in association for the advancement of beliefs and ideas is an inseparable aspect of the `liberty' assured by the Due Process Clause of the Fourteenth Amendment, which embraces freedom of speech. \dots\ [S]tate action which may have the effect of curtailing the freedom to associate is subject to the closest scrutiny.''
\end{quote}
\noindent --- Justice Harlan, 357 U.S. at 460--461.

\end{document}
